%============================================================
\chapter{情報構造の供給者：手続きによる決済}
\label{ch:protocol}
%============================================================

\section{本章で課す前提}

第\ref{ch:solo}章は資源に制約を課し、第\ref{part:unmanned}章はその極限をとった。
いずれも第\ref{sec:info}章の情報構造 $\mathcal{G}$ を所与としている。
本章は $\mathcal{G}$ 自体を動かす。すなわち\textbf{誰が $\mathcal{G}$ を供給するかを差し替えたとき、
実現可能な $\Phi$ の集合がどう変わるか}を問う。

二点を先に断る。

第一に、\textbf{本章は実証ではなく演繹である。} 結論の妥当性は前提の妥当性に依存する。

第二に、\textbf{本章の内容に新規性はない。} ここで述べる関係はいずれも既存文献に対応があり、
付録\ref{app:priorwork}に対応表を置いた。
本章の目的は、第\ref{part:theory}部で構成した理論が
具体的な条件の下でどう機能するかを示すことにある。

\begin{definition}[手続き]
\label{def:protocol}
事前に公開され、実行結果を全参加者が独立に検証できる規則の集合を\emph{手続き}と呼び、
$\mathcal{P}$ と書く。$\mathcal{P}$ が保持する変数の組を $\mathcal{P}$ の\emph{状態}という。
状態遷移は不可分であるとする。すなわち遷移の一部だけが適用されることはない。
\end{definition}

前提は次のとおりである。

\begin{enumerate}[label=(\arabic*)]
\item 決済 $\pi$ は $\mathcal{P}$ 上の状態遷移として表現される
\item $\mathcal{P}$ の実行結果は全参加者が独立に検証できる（定義\ref{def:protocol}）
\item 履行 $\delta$ が $\mathcal{P}$ の状態に含まれるかは対象による
\end{enumerate}

前提(1)(2)を満たす実装として現に広く用いられているのは、
公開型の分散台帳、すなわちブロックチェーン上の決済である。
ただしそう呼ばれる仕組みのすべてが前提を満たすわけではない。
検証者が限定される場合、および状態を特定の事業者が保持する場合は(2)を満たさず、
本章の演繹は適用できない。

(1)と(2)は緩和である。第\ref{sec:info}章の $\mathcal{G}=\bigcap_i\mathcal{G}_i$ を供給するのは
通常は当事者の外にある第三者であり、$\pi$ の検証と執行はそこに依存していた。
前提(1)(2)はこれを不要にする。

\textbf{(3)のみ前提として固定できない。} これが本章の場合分けを要求する。

\begin{center}
\begin{tabularx}{\textwidth}{@{}llX@{}}
\toprule
\multicolumn{3}{@{}l}{\textbf{本章で扱う理論量}}\\
\midrule
$\mathcal{G}$ & 第\ref{sec:info}章 & 供給者を差し替える対象 \\
$\kappa_i$ の添字 $i$ & 式\eqref{eq:kappa} & 手続き上では識別子になる。\textbf{加法性が失われる} \\
$\dbar-\delta$ & 定義\ref{def:breakage-def} & 区分「内」では手続き上の状態として観測できる \\
$E$ & 式\eqref{eq:cash} & 担保として現れる \\
自由度(2)(3) & 第\ref{sec:dof}節 & 篩がこの二つで切れる \\
\bottomrule
\end{tabularx}
\captionof{table}{本章で扱う理論量。情報構造の供給者を差し替えたときに関わる記号。}
\label{tab:protocol-quantities}
\end{center}

\section{前提の翻訳}

\subsection{第三者が担っていた三つの機能}

$\mathcal{G}$ の供給者が契約において果たしている役割を三つに分ける。

\begin{center}
\begin{tabularx}{\textwidth}{@{}lXl@{}}
\toprule
機能 & 内容 & 通常の担い手 \\
\midrule
(a) 記録 & $\kappa_i$ の状態を保持する & 銀行、決済代行 \\
(b) 検証 & $\pi$ の発生条件が成立したかを判定する & 裁判所、監査 \\
(c) 執行 & 条件成立時に価値を移転する & 裁判所、銀行 \\
\bottomrule
\end{tabularx}
\captionof{table}{第三者が契約において果たす三つの機能。}
\label{tab:third-party-functions}
\end{center}

前提(1)(2)より、$\mathcal{P}$ は(a)を常に担える。
(b)は条件が $\mathcal{P}$ の状態から判定できる限り担える。
(c)は移転の対象が $\mathcal{P}$ の状態である限り担える。

\subsection{履行の所在による場合分け}

前提(3)により二つの場合が生じる。

\begin{center}
\begin{tabularx}{\textwidth}{@{}llXX@{}}
\toprule
区分 & $\delta$ の所在 & $\mathcal{G}$ に載るもの & $\mathcal{P}$ が担える機能 \\
\midrule
内 & $\mathcal{P}$ の状態 & $\delta$ と $\pi$ の両方 & (a)(b)(c) すべて \\
外 & $\mathcal{P}$ の外 & $\pi$ のみ & (a)、および $\pi$ に関する(b)(c) \\
\bottomrule
\end{tabularx}
\captionof{table}{履行の所在による区分と、手続きが担える機能。}
\label{tab:protocol-cases}
\end{center}

区分の差は第\ref{sec:dof}節の自由度(1)、すなわち時間関係に直接現れる。

\begin{proposition}[同時性]
\label{prop:simultaneity}
区分「内」では、任意の $\Phi$ に対し $\kappa_i=0$ となる実装が存在する。
区分「外」では、$\delta$ を伴う $\Phi$ について $\kappa_i=0$ を実現できない。
\end{proposition}

\begin{proof}
区分「内」では $\delta$ と $\pi$ がともに $\mathcal{P}$ の状態遷移である。
定義\ref{def:protocol}の不可分性より両者を単一の遷移として構成でき、
このとき $D_i$ と $P_i$ が同時に動くから式\eqref{eq:kappa}の $\kappa_i$ は恒等的に $0$ となる。

区分「外」では $\delta$ が $\mathcal{P}$ の状態でないため、
$\mathcal{P}$ は $\delta$ の時刻を持たない。
$\kappa_i=0$ は $\tau_i=0$ を要求するが、$\tau_i$ は $\mathcal{G}$ 上で確定しない。
\end{proof}

すなわち\textbf{区分「内」では自由度(1)が完全な設計変数になり、
区分「外」では $\kappa\neq 0$ が強制される}。
現実の即時交換は当事者が互いを同時に観測することで同時性を作っているが、
$\mathcal{P}$ はその観測を持たない。
第\ref{sec:protocol-filter}節で 1-1（即時交換）が区分「外」で消えるのはこのためである。

\begin{remark}[区分は対象の属性ではなく $\Phi$ の属性である]
\label{rem:case-is-phi}
同一の対象であっても、何を履行とみなすかにより区分が変わる。
手続き上の残高を他の残高と交換する $\Phi$ では区分は内であり、
同じ残高を手続きの外にある通貨と交換する $\Phi$ では外になる。

所有の記録が手続き上にあることは、区分を内にしない。
非代替性トークン（NFT）を現物の美術品や会員権に紐づける形式では、
移転するのは記録であり、履行は現実の引渡しである。
$\delta$ が $\mathcal{P}$ の外にあるから区分は外になる。
両者を同一視できないことは、記録の保有者が変わっても引渡しが起きないという形で現れる。
同じトークンであっても、手続き上の権利の行使のみを履行とする $\Phi$ では区分は内になる。

したがって「その対象は内か外か」という問いは立たない。
問えるのは「その $\Phi$ は内か外か」である。
第\ref{sec:unit}節の結論と同じ形であり、観測単位を対象ではなく $\Phi$ に置くことがここでも要求される。
\end{remark}

\begin{remark}[区分「外」における仲介]
\label{rem:oracle}
区分「外」において $\delta$ に関する条件を $\mathcal{P}$ に持ち込むには、
外の事実を状態として書き込む主体が要る。
$\mathcal{P}$ が検証できるのは所定の署名が付いているかまでであり、
書き込まれた値が事実かは検証していない。
すなわちこの主体は前提(2)を満たさず、\textbf{除去したはずの第三者が戻る}。

書き込みの主体は単一の場合も複数の合議による場合もあり、
除去の程度は離散的でなく連続的である。
本章はこの連続体を扱わず、両端のみを扱う（第\ref{sec:protocol-limits}節）。
実装の分類は \cite{pasdar2023} にある。
\end{remark}

\subsection{拘束する制約の入れ替わり}

第\ref{ch:solo}章では、財務制約と容量制約のいずれが先に拘束するかが
前提により入れ替わった。本章でも同じことが起きる。

前提(1)(2)は $\mathcal{G}$ を広げるため、命題\ref{prop:implementable}の可測性条件は
区分「内」において本稿のどの領域よりも緩む。
\textbf{したがって $\mathcal{G}$ は拘束しなくなり、別の量が拘束する。}
それが何であるかを第\ref{sec:collateral}節で特定する。

\section{識別子と主体の乖離}
\label{sec:identifier}

本節は副産物である。第\ref{ch:solo}章がクラウドについて
$A\to\kappa_S<0$ の対応を記録したのと同じ位置に置く。

式\eqref{eq:kappa}の $\kappa_i$ は主体 $i$ ごとに定義され、
式\eqref{eq:cash}では $\sum_i\kappa_i$ として加算される。
この加算は\textbf{添字が主体を一意に指すこと}を前提としている。

手続き上で同定できるのは識別子である。
一つの主体が複数の識別子を持つことを $\mathcal{P}$ は排除できない。
識別子の生成に費用がかからないためであり、これは既知の問題である \cite{douceur2002}。

\begin{remark}[加法性の喪失]
\label{rem:additivity}
識別子を $j$、主体を $i$、対応を $j\mapsto\iota(j)$ とすると、
観測できるのは $\kappa_j$ であり、必要なのは
\begin{equation}
\kappa_i = \sum_{j:\,\iota(j)=i} \kappa_j
\label{eq:kappa-aggregation}
\end{equation}
である。$\iota$ が観測できないため、式\eqref{eq:kappa-aggregation}は評価できない。

\textbf{これは第\ref{sec:phi}章の定義の限界であって、本章の前提の帰結ではない。}
$N$ を固定し添字を主体に固定した設定が、識別子が可変な環境に合わないことを示している。
第\ref{sec:token-limit}節で述べた請求権の流通と同じ箇所、
すなわち定義域の設定に由来する。理論の修正は行わない。
\end{remark}

注\ref{rem:additivity}は族7に直接効く。

\begin{proposition}[受益者の同定を要する $\Phi$ は実装できない]
\label{prop:beneficiary}
$\iota$ が観測できないとき、受益者と支払者が異なることを条件とする $\pi$ は
$\mathcal{G}$ 可測でない。
\end{proposition}

\begin{proof}
族7の定義的性質は、$\delta$ を受ける主体と $\pi$ を出す主体が異なることである。
$\mathcal{P}$ が同定するのは識別子 $j$ であり $\iota(j)$ は観測できないから、
$\iota(j_1)\neq\iota(j_2)$ は $\mathcal{G}$ 上で判定できない。
命題\ref{prop:implementable}より、この条件に依存する $\pi$ は書けない。
\end{proof}

命題\ref{prop:beneficiary}は区分に依存しない。
$\delta$ が $\mathcal{P}$ の状態にあっても、それを受ける\emph{主体}は同定できないためである。
族7のうち 7-1 と 7-2 が両区分で脱落する理由はここにある。
7-4 が残るのは、対価関係を明示せず受益者の同定を要しないためである。

識別子の生成が自由であるとき、主体の相異を前提とする仕組みが機能しないことは、
メカニズムデザインでは false-name-proofness の問題として扱われている
\cite{yokoo2004}。命題\ref{prop:beneficiary}はその契約形態への現れである。

\section{類型のふるい分け}
\label{sec:protocol-filter}

第\ref{part:catalog}部の28類型を、表\ref{tab:protocol-cases}の二区分で篩にかける。

\textbf{篩の基準は第\ref{ch:solo}章とは異なる。} 第\ref{ch:solo}章は実行可能かどうかで切ったが、
ここでは手続きが寄与するかどうかで切る。区分「外」で寄与を受けない類型も、
執行が第三者に戻るだけで消滅はしない。
以下、$\pi$ を規定する条件が $\mathcal{G}$ に載り、
かつ執行の対象が $\mathcal{P}$ の状態にあるものを「寄与を受ける」とする。
判定は類型ごとではなく $\Phi$ ごとに行う（注\ref{rem:case-is-phi}）。

% 28類型のふるい分け。判定の行は下の二つの macro にしか実体がない。
% PDF では紙面の高さに収まらないため族1--4と族5--7の二表に分けるが、
% HTML にはその制約がないので一つにまとめる。
% 判定を直すときは macro だけを見ればよい。
\newcommand{\protocolfilterrowsA}{%
\multicolumn{5}{@{}l}{\textbf{族1：単発交換}}\\
1-1 & 即時交換 & ○ & --- & 同時性は $\delta$ と $\pi$ が同一の状態遷移である場合にのみ作れる \\
1-2 & 前受単発 & ○ & ○ & $\pi$ が先行し、$\delta$ を参照しない \\
1-3 & 後払単発 & ○ & --- & 外では $\delta$ の完了を判定できない \\
1-4 & 分割後払 & ○ & 部分 & 外でも $\pi$ の徴収は担えるが、不履行時に原資産へ及ばない \\
\midrule
\multicolumn{5}{@{}l}{\textbf{族2：継続アクセス}}\\
2-1 & 定額アクセス & ○ & ○ & $\pi=c$ が $\delta$ を参照しない \\
2-2 & 従量 & ○ & --- & 外では $\delta$ の計量に注\ref{rem:oracle}の主体を要する \\
2-3 & 二部料金 & ○ & 部分 & 外では固定部分 $F$ のみ \\
2-4 & プリペイド & ○ & ○ & 内では未使用残高が手続き上の状態そのものになる \\
2-5 & オプション・保証 & ○ & --- & 外では行使条件が手続きの外にある \\
\midrule
\multicolumn{5}{@{}l}{\textbf{族3：資産の時間貸し}}\\
3-1 & リース・レンタル & ○ & --- & 内は資産自体が手続き上の状態である場合に限る \\
3-2 & 共用資産 & 条件付 & --- & 稼働を状態遷移として書ける範囲が狭い \\
3-3 & 譲渡＋長期回収 & ○ & --- & 内は担保付き。外では原資産に及ばない \\
\midrule
\multicolumn{5}{@{}l}{\textbf{族4：補完財による回収}}\\
4-1 & ロックイン消耗品 & 条件付 & --- & 本体が手続きの外にあるという族の前提が内では崩れる \\
4-2 & フリーミアム & ○ & ○ & $\kappa>0$ だが自ら発行した信用であり、執行を手続きに求めていない \\
4-3 & 本体補填＋金融 & ○ & --- & 外では本体が手続きの外にある \\
}
\newcommand{\protocolfilterrowsB}{%
\multicolumn{5}{@{}l}{\textbf{族5：成果・状態連動}}\\
5-1 & 成果報酬 & 条件付 & --- & $\pi=h(y)$ の $y$ が手続き上の量である場合に限る \\
5-2 & レベニューシェア & ○ & --- & 内では相手方の流量が手続き上で観測できる \\
5-3 & コストプラス & --- & --- & 実費が内外いずれでも手続きの外にある \\
5-4 & 保険引受 & ○ & --- & 内は $\omega$ が手続き上の事象に限られる \\
5-5 & IPライセンス & ○ & --- & 内は利用が手続き上で観測できる場合 \\
\midrule
\multicolumn{5}{@{}l}{\textbf{族6：媒介}}\\
6-1 & 取引額手数料 & ○ & --- & 内では 1-1 の同時性と組み合わせて完結する \\
6-2 & エスクロー預り & ○ & 部分 & 外では片側の預りのみ。解放条件に注\ref{rem:oracle}の主体を要する \\
6-3 & スプレッド & ○ & --- & 内では在庫が手続き上の残高になる \\
6-4 & 機会課金 & 条件付 & --- & 接触を手続き上の事象として定義できる範囲に限る \\
\midrule
\multicolumn{5}{@{}l}{\textbf{族7：受益者と支払者の分離}}\\
7-1 & 広告 & --- & --- & 受益者と支払者が異なることを検証できない（第\ref{sec:identifier}節） \\
7-2 & クロスサブシディ & --- & --- & 同上 \\
7-3 & 公的支払 & --- & --- & 支払者が価格を規定し、設計自由度が制度層にある \\
7-4 & 寄付・支援 & ○ & ○ & 対価関係を明示しないため受益者の同定を要しない \\
}

\ifdefined\HCode
\begin{center}
\begin{tabularx}{\textwidth}{@{}llccX@{}}
\toprule
番号 & 類型 & 内 & 外 & 判定の根拠 \\
\midrule
\protocolfilterrowsA
\midrule
\protocolfilterrowsB
\bottomrule
\end{tabularx}
\captionof{table}{28類型のふるい分け。○は寄与を受けるもの。}
\label{tab:protocol-filter-a}
\end{center}
\else
\begin{center}
\begin{tabularx}{\textwidth}{@{}llccX@{}}
\toprule
番号 & 類型 & 内 & 外 & 判定の根拠 \\
\midrule
\protocolfilterrowsA
\bottomrule
\end{tabularx}
\captionof{table}{28類型のふるい分け（族1--4）。○は寄与を受けるもの。}
\label{tab:protocol-filter-a}
\end{center}

\begin{center}
\begin{tabularx}{\textwidth}{@{}llccX@{}}
\toprule
番号 & 類型 & 内 & 外 & 判定の根拠 \\
\midrule
\protocolfilterrowsB
\bottomrule
\end{tabularx}
\captionof{table}{28類型のふるい分け（族5--7）。○は寄与を受けるもの。}
\label{tab:protocol-filter-b}
\end{center}
\fi

区分「内」で寄与を受けるのは20類型、区分「外」では5類型である。
第\ref{ch:solo}章の篩が28類型を10前後に減らしたのに対し、
区分「内」の篩はほとんど減らさない。
\textbf{本章の前提は制約ではなく緩和だからである。}

区分「外」で 1-1 が消えるのは命題\ref{prop:simultaneity}による。
$\delta$ が手続きの外にあれば前受か後払のいずれかに退化する。

\textbf{ただし、寄与を受けることと実行できることは別である。}
前受型（1-2、2-4）は表の上では区分「内」で通るが、
実行には相手方に対する担保を要し、資産を持たない主体には開かれていない
（系\ref{cor:growth-needs-capital}、系\ref{cor:exclusion}）。
第\ref{sec:collateral}節で扱う。

\begin{remark}[篩が切れる軸]
区分「外」で残る5類型は 1-2、2-1、2-4、4-2、7-4 であり、
共通するのは $\pi$ が $\delta$ に可測でないことのみである。
第\ref{sec:dof}節の六自由度のうち、自由度(2)だけがこの篩に情報を持つ。

\textbf{これは命題ではない。} 区分「外」を $\delta$ が $\mathcal{G}$ に載らない場合として
定義したのだから、$\delta$ 可測な $\pi$ が書けないことは定義からの帰結である。
情報があるのは主張の側ではなく、28類型を個別に当てた結果が
単一の自由度で説明されたという事実の側にある。
\end{remark}

\begin{remark}[制度層の作用の向き]
2-4（プリペイド）と 6-2（エスクロー）は、第\ref{sec:bootstrap}節では
供託義務と登録制により脱落した類型であった。
区分「内」では手続きとして実装される。

\textbf{同一の類型が、制度層の作用の向きにより脱落もし復帰もする。}
注\ref{rem:institution-modes}の三形式のうち第一の形式（参入要件）について、
作用が一方向でないことを示す例になる。
\end{remark}

\section{担保制約}
\label{sec:collateral}

第\ref{sec:protocol-filter}節で $\mathcal{G}$ が拘束しなくなることを見た。
本節で何が拘束するかを特定する。

区分「内」において $\kappa_i>0$、すなわち事業者が与信を出す $\Phi$ を考える。
$\pi$ が不履行となったとき執行するには、執行の対象が $\mathcal{P}$ の状態になければならない。
一方、$\mathcal{G}$ は $\mathcal{P}$ の状態のみからなるため、
\textbf{相手方の弁済能力も過去の履行実績も $\mathcal{G}$ に含まれない}。

したがって与信は事前に拘束された状態、すなわち担保によってのみ手当てできる。
自己が相手方 $i$ のために拘束した担保を $B_i$、
相手方 $i$ が自己のために拘束した担保を $B_i'$ と書くと、

\begin{equation}
\max\{\kappa_i,0\} \;\leq\; B_i',
\qquad
\max\{-\kappa_i,0\} \;\leq\; B_i
\label{eq:protocol-collateral}
\end{equation}

が成立する。担保は自己の資産から重複なく供されるから $\sum_i B_i\leq E$ であり、次を得る。

\begin{proposition}[担保制約]
\label{prop:collateral}
区分「内」において、受けられる与信の総額は自己資本を超えない。
\begin{equation}
\sum_i \max\{-\kappa_i,0\} \;\leq\; E
\label{eq:protocol-finance}
\end{equation}
\end{proposition}

\begin{proof}
式\eqref{eq:protocol-collateral}の第二式について全ての $i$ の和をとると
$\sum_i\max\{-\kappa_i,0\}\leq\sum_i B_i$ を得る。
同一の資産を複数の相手方に対して同時に拘束することはできないから
$\sum_i B_i\leq E$ であり、両者を合わせる。
\end{proof}

式\eqref{eq:protocol-finance}を第\ref{ch:solo}章の式\eqref{eq:solo-finance}と比較する。
そこでは $E\approx 0$ の下で $\sum\max\{-\kappa_i,0\}\geq\sum\max\{\kappa_i,0\}$、
すなわち引き出した与信が与えた与信を上回ることが要求された。
本章では引き出せる与信そのものが $E$ で抑えられる。

\begin{corollary}[代替が片方向に退化する]
\label{cor:one-way-substitution}
区分「内」において、命題\ref{prop:substitution}の代替関係は片方向になる。
資本は信用を代替するが、信用は資本を代替しない。
\end{corollary}

\begin{proof}
命題\ref{prop:collateral}より $\kappa<0$ の項は $E$ を上限とする。
式\eqref{eq:cash}において $\kappa<0$ の項が $E$ を超えて調達を増やすことはない。
逆向きの代替、すなわち $E$ が $\kappa<0$ を代替することは
命題\ref{prop:substitution}のとおり成立する。
\end{proof}

系\ref{cor:one-way-substitution}は、命題\ref{prop:substitution}が
本稿で三度目に現れる箇所である。

\begin{center}
\begin{tabularx}{\textwidth}{@{}lllX@{}}
\toprule
箇所 & $E$ & 信用 & そこで拘束するもの \\
\midrule
第\ref{ch:solo}章 一人事業 & 0 & 唯一の調達手段 & 信用のブートストラップ問題 \\
第\ref{sec:retiree}節 資産を持つ層 & 正 & ゼロから積む & 資本が信用を積む期間を買う \\
本章 区分「内」 & 唯一の調達手段 & 0 & 積む期間が存在しない \\
\bottomrule
\end{tabularx}
\captionof{table}{命題\ref{prop:substitution}が現れる三つの箇所。両端が揃う。}
\label{tab:substitution-three-points}
\end{center}

第\ref{sec:credit-formation}節は、公開された履行の蓄積により信用が形成される経路を述べた。
区分「内」ではこの経路が働かない。
$\mathcal{P}$ が識別子しか同定できず（第\ref{sec:identifier}節）、
過去の履行実績を主体に帰属させられないためである。

\begin{corollary}[資産を持たない主体の排除]
\label{cor:exclusion}
$A=\mathrm{Inv}=0$ かつ $E=0$ のとき、実現可能なのは
全ての $i$ について $\kappa_i=0$ である $\Phi$ に限られる。
\end{corollary}

\begin{proof}
式\eqref{eq:protocol-finance}に $E=0$ を代入すると
全ての $i$ について $\kappa_i\geq 0$ となる。
一方、命題\ref{prop:substitution}より $A=\mathrm{Inv}=E=0$ のとき
$\sum_i\max\{-\kappa_i,0\}\geq\sum_i\max\{\kappa_i,0\}$ が要求され、
左辺は $0$ であるから $\kappa_i\leq 0$ を得る。
\end{proof}

\begin{corollary}[前受による成長は資本を要する]
\label{cor:growth-needs-capital}
区分「内」において $\CCC<0$ を実現するには $E>0$ を要する。
\end{corollary}

\begin{proof}
式\eqref{eq:ccc}より $\CCC=W/r$、$W=\sum_i\kappa_i+\mathrm{Inv}$ である。
命題\ref{prop:collateral}に $E=0$ を代入すると全ての $i$ について $\kappa_i\geq 0$ となり、
$\mathrm{Inv}\geq 0$ と併せて $W\geq 0$、したがって $\CCC\geq 0$ を得る。
\end{proof}

第\ref{ch:solo}章は、無資産の一人事業に対して年払いの前受により
$\CCC<0$ を作る経路を示した。成長が現金を生む状態である。
\textbf{この経路は区分「内」では閉じる。}
前受を受けるには相手方に対する担保が要り、担保は資本から供されるためである。

すなわち\textbf{資産を持たない主体に残るのは $\kappa$ が恒等的に $0$ の類型に限られる}。
第\ref{part:catalog}部で $\kappa\approx 0$ とした 1-1（即時交換）、2-2（従量）、
6-1（取引額手数料）とその合成がこれにあたる。
第\ref{ch:solo}章では $E=0$ が $\kappa_C<0$ を要求したが、
ここでは $\kappa=0$ を強制する。
第\ref{sec:bootstrap}節の解(a)(b)(c)はいずれも信用の不足を前提としていたが、
ここで不足するのは資本である。

\begin{remark}[既存理論との対応]
\label{rem:collateral-priorwork}
本節の内容は新規ではない。純資産が担保能力を通じて調達を制約する構造は
\cite{kiyotaki1997}に、担保制約が企業の資金調達と危険管理を同時に規定する定式化は
\cite{rampini2010,rampini2013}にある。
命題\ref{prop:substitution}自体もこれらに対応がある（付録\ref{app:priorwork}）。

前提(1)(2)の下で契約可能な空間が広がることは\cite{cong2019}が扱っている。
系\ref{cor:exclusion}に相当する指摘、すなわち担保を要求する仕組みが
資本を持たない主体を排除するという論点も、既に文献がある\cite{oecd2024}。

\textbf{本節の役割は、これらを式\eqref{eq:cash}の同一の不等式の上に置くことにある。}
\end{remark}

\section{観測可能性}
\label{sec:protocol-observability}

前提(2)は副次的な帰結を持つ。
$\mathcal{P}$ の状態が全参加者に対して検証可能であることから、
区分「内」では次の量が個別に観測できる。

\begin{center}
\begin{tabularx}{\textwidth}{@{}lXX@{}}
\toprule
量 & 本稿における測定状況 & 区分「内」での状態 \\
\midrule
$M(t)$ & 個社データを要し、集計では意味を持たない & 識別子ごとに観測できる \\
$\kappa_j$ & 同上 & 同上 \\
$\dbar-\delta$ & 集計値からの推定にとどまる & 未使用残高として直接観測できる \\
$B$ & 対応する量を本稿は持たない & 拘束された状態として観測できる \\
\bottomrule
\end{tabularx}
\captionof{table}{区分「内」において観測可能になる量。}
\label{tab:protocol-observability}
\end{center}

ただし観測できるのは識別子単位であり、注\ref{rem:additivity}により
主体単位の $\kappa_i$ は復元できない。
\textbf{測定の障害が、集計の不在から集計の不能へ入れ替わる。}
第\ref{part:method}部の障害分類に対して、これは新たな形式にあたる。

\subsection{母集団フレームが得られる}
\label{sec:protocol-frame}

観測可能になるのは量だけではない。

第\ref{ch:method}章の生存条件づけは、
観測される対象が制約\eqref{eq:constraint}を破らなかった経路に限られることから生じる。
区分「内」では\textbf{手続きの記録そのものが母集団になる}。
稼働した $\Phi$ は停止したものも含めて残るため、条件づけが発生しない。

第\ref{ch:platformfee}章は無人継続期間の直接測定を
「母集団フレームを欠く」として断念した。断念の理由がここでは消える。

\subsection{検証できることとできないこと}
\label{sec:protocol-testable}

観測できることは検証できることを意味しない。三点を先に区別する。

\begin{center}
\begin{tabularx}{\textwidth}{@{}lX@{}}
\toprule
区別 & 内容 \\
\midrule
手続きが強制する量 & 観測しても検証にならない。規約を読んでいるにすぎない \\
設計の結果である量 & 検証の対象になる \\
主体単位の量 & 復元できない（注\ref{rem:additivity}） \\
\bottomrule
\end{tabularx}
\captionof{table}{区分「内」で観測できる量の三区分。}
\label{tab:protocol-testability}
\end{center}

命題\ref{prop:collateral}の担保制約は第一に属する。
手続きが担保を要求する以上、担保が置かれていることを観測しても命題の検証にならない。
\textbf{検証の対象になるのは、事業者が選んだ結果として現れる量に限られる。}

\begin{remark}[母集団が異なる]
\label{rem:protocol-population}
区分「内」で測った分布を企業一般に一般化することはできない。
第\ref{ch:method}章の意味で別の母集団である。
\textbf{ここで検証できるのは演繹が正しいかであって、
その演繹が現実の一人事業に当てはまるかではない。}
第\ref{ch:solo}章の命題は検証できるが、同章の対象については何も言わない。
\end{remark}

\subsection{予測の登録}
\label{sec:protocol-predictions}

第\ref{ch:next}章の手続きに従い、取得の前に予測を固定する。
本章の予測は四つあるため、$\Phi$ に番号を付して区別する。

\begin{quote}
\textbf{予測$\Phi$1}：手続きが呼び出す他の手続きの数が少ないほど、
改修されるまでの期間が長い。命題\ref{prop:unmanned-period}の $\lambda_{\rm dep}$ 側の含意である。

\textbf{対抗仮説}：改修の頻度は依存の数ではなく利用量で決まる。
依存が少なくとも利用が多ければ改修される。
\end{quote}

\begin{quote}
\textbf{予測$\Phi$2}：取引が途絶えた $\Phi$ について、
$M(t)$ は途絶に先行して単調に減少する。
第\ref{sec:objective}章の $\tau=\inf\{t:M(t)<0\}$ が吸収状態であることの含意である。

\textbf{対抗仮説}：途絶は $M(t)$ と無関係に生じる。
$M(t)$ は途絶の直前まで変化しない。
\end{quote}

\begin{quote}
\textbf{予測$\Phi$3}：区分「内」で稼働する $\Phi$ の $\CCC$ の分布は非負に偏る。
系\ref{cor:growth-needs-capital}より、$\CCC<0$ の実現には $E>0$ を要するためである。

\textbf{対抗仮説}：$\CCC$ の分布は第\ref{part:industry}章の企業部門と同様であり、
区分「内」に固有の偏りはない。
\end{quote}

\begin{quote}
\textbf{予測$\Phi$4}：稼働する $\Phi$ を第\ref{sec:type-assignment}節の規則で割り当てると、
族7は 7-4 を除いてほぼ現れない。命題\ref{prop:beneficiary}の含意である。

\textbf{対抗仮説}：族7型も他の族と同程度の比重で現れる。
\end{quote}

\paragraph{事前に確定した事項。}

\begin{center}
\begin{tabularx}{\textwidth}{@{}lX@{}}
\toprule
項目 & 内容 \\
\midrule
母集団 & 基準時点に手続きの記録上で稼働していた $\Phi$ の全数。事後に選ばない \\
下限 & 取引件数が基準期間に一定数未満のものを除く。閾値は取得前に確定する \\
途絶の定義 & 一定期間にわたり取引が生じないこと。区分「内」では手続き自体は停止しないため、停止ではなく途絶で定義する \\
定常状態 & $\CCC$ は定常を前提とするため、$r$ が安定した期間に限る（第\ref{sec:age}節） \\
観測単位 & 識別子。主体単位への集計は行わない（注\ref{rem:additivity}） \\
類型の割り当て & 第\ref{sec:type-assignment}節の判定順序を用いる。順序は変更しない \\
\bottomrule
\end{tabularx}
\captionof{table}{予測$\Phi$1--$\Phi$4 について事前に確定した事項。}
\label{tab:protocol-registration}
\end{center}

\begin{remark}[$\lambda_{\rm dep}$ と利用量は相関しうる]
\label{rem:dep-usage-confound}
予測$\Phi$1 の対抗仮説は棄却しにくい。
依存の多い手続きは機能が豊富であり、利用量も大きい傾向があるためである。
検定では利用量を統制する必要があるが、
統制後も残差に依存の効果が現れるかは事前には分からない。
\textbf{統制の手続きを取得前に確定しておく。}
\end{remark}

\subsection{予測の検定}
\label{sec:protocol-results}

公開型の分散台帳の記録により四つを検定した。
母集団は基準時点で稼働していた $\Phi$ の全数であり、標本は事後に選んでいない。
支持されたのは一つである。

\begin{center}
\begin{tabularx}{\textwidth}{@{}llXl@{}}
\toprule
予測 & 標本 & 結果 & 判定 \\
\midrule
$\Phi$1 & 改修可能な手続き 193 万件 & 依存の数は改修までの期間を説明しない & 棄却 \\
$\Phi$2 & 手数料と残高の双方を持つ 427 件 & 途絶前 30 日で減少（$p=3.2\times10^{-9}$） & 支持 \\
$\Phi$3 & 稼働する $\Phi$ 52 件 & 対抗仮説が識別されない & 検証不能 \\
$\Phi$4 & 同上 & 族7 は 1 件だが 7-2 であり 7-4 は現れない & 棄却 \\
\bottomrule
\end{tabularx}
\captionof{table}{予測$\Phi$1--$\Phi$4 の検定結果。}
\label{tab:protocol-results}
\end{center}

\begin{remark}[$\Phi$1：複製の集中が推定を支配する]
\label{rem:clone-concentration}
193 万の識別子が指す実装は 76,609 種しかない。
うち一つを 155 万件、すなわち母集団の 80\% が指している。
これらは独立な $\Phi$ ではなく、一つの $\Phi$ の複製である。

全体では依存の数の係数が正で有意になるが（$p=0.004$）、
この一つのクラスタを外すと符号が逆転して有意でなくなる（$p=0.45$）。
\textbf{クラスタ頑健標準誤差では防げない。}
クラスタ内の相関は補正するが、一つのクラスタが曝露の 95\% を占める場合の
点推定の偏りは補正しないためである。

注\ref{rem:additivity}の加法性の喪失は主体と識別子の乖離を述べたが、
同一の実装を指す識別子の集中は別の形の乖離である。
\textbf{台帳上で標本を組むときは、複製の集中を先に数える必要がある。}
\end{remark}

\begin{remark}[$\Phi$3：対照が識別されない]
\label{rem:phi3-illposed}
$\CCC<0$ は前受けであり、系\ref{cor:growth-needs-capital}によれば $E>0$ を要する。
しかし比較対象である企業部門の $\CCC$ は、
第\ref{part:industry}章の 58 業種すべてで非負である（最小 4.9 日）。
負の $\CCC$ は個別企業に現れるもので、業種の集計では相殺されて消える。

予測と対抗仮説が同じことを述べており、観測しても区別がつかない。
第\ref{sec:predB-illposed}節の予測Bと同型の不良設定である。
加えて区分「内」の $\tau$ は秒の単位であり、企業部門の日の単位とは桁が五つ違う。
\end{remark}

\begin{remark}[$\Phi$4：実装できることと現れることは別である]
\label{rem:implementable-not-observed}
族7 は 52 件中 1 件であり、稀であるという量的な部分は当たった。
しかし現れたのは 7-2（クロスサブシディ）であり、7-4 は 0 件であった。

予測は命題\ref{prop:beneficiary}から「7-4 のみが残る」を導いていたが、
命題が述べるのは実装の可否であって出現の頻度ではない。
\textbf{実装できることから現れることを導いたのが誤りである。}

なお唯一の 7-2 は、受益者の選定を署名により手続きの外で行っている。
命題\ref{prop:beneficiary}と整合するが、
予測の文言が外れている以上これを支持の証拠として扱わない。
\end{remark}

\begin{remark}[判定順序の前に一段が要る]
\label{rem:is-it-phi}
表\ref{tab:family-assignment-order}は $\Phi$ を所与としている。
台帳の識別子に当てると、標本 200 件のうち 131 件が $\Phi$ ではなかった。
トークンの台帳（決済の媒体であり $\delta$ と $\pi$ の対を持たない）が 74 件、
対価を伴わない基盤が 27 件、
そして\textbf{決済媒体の名を模した識別子が 30 件}である。

最後のものは定義\ref{def:protocol}を満たすが履行を持たない。
第\ref{sec:identifier}節の $\iota$ の観測不能が、
名称の詐称という形で現れている。
\end{remark}

\subsection{副次的に測れた量}
\label{sec:protocol-byproducts}

検定の過程で、本稿が他の章で測定を断念していた量が測れた。

\paragraph{未想定状態への応答の率。}
命題\ref{prop:unmanned-period}は無人で運転できる期間を $1/\lambda_{\rm novel}$ とした。
$\lambda_{\rm novel}$ は本稿のどこでも測定されていない。

区分「内」では二つの率が測れる。

\begin{center}
\begin{tabularx}{\textwidth}{@{}lXll@{}}
\toprule
量 & 内容 & $\lambda$（毎年） & $1/\lambda$ \\
\midrule
改修の率 $\lambda_{\rm fix}$ & 実装が差し替えられた率。193 万件、うち複製の集中を除く & 0.048 & 20.6 年 \\
停止の率 $\lambda_{\rm stop}$ & 稼働が絶えた率。1,718 件の全数 & 0.030 & 33.8 年 \\
\bottomrule
\end{tabularx}
\captionof{table}{区分「内」で測った応答の率。}
\label{tab:lambda-measured}
\end{center}

\begin{remark}[測っているのは $\lambda_{\rm novel}$ ではない]
\label{rem:lambda-is-response}
表\ref{tab:lambda-measured}の二つは、いずれも\textbf{未想定状態の到来}ではなく
\textbf{到来に対する応答}を数えている。改修は運営者が応答した場合であり、
停止は応答しなかったか、応答できなかった場合である。

到来しても応答を要さなかった場合は、どちらにも現れない。したがって
\[
\lambda_{\rm novel}\ \geq\ \lambda_{\rm fix}+\lambda_{\rm stop}
\]
であり、測れたのは下限である。母集団が異なるため二つは加法的でないが、
\textbf{命題\ref{prop:unmanned-period}の $1/\lambda_{\rm novel}$ は 20 年より短い}という上限が得られる。

改修までの期間の分布は極端に偏る。改修された 20,732 件の中央値は 42 日であるのに対し、
改修されないまま観測を終えた 190 万件の中央値は 4.6 年である。
\textbf{手を入れるものは直ちに入れ、入れないものは何年も入れない。}
平均としての $1/\lambda$ はこの二峰性を覆い隠す。
\end{remark}

\paragraph{手続きの継続期間。}
第\ref{ch:platformfee}章は無人継続期間の直接測定を「母集団フレームを欠く」として断念した。
第\ref{sec:protocol-frame}節のとおり、区分「内」ではこの断念の理由が消える。

稼働した 1,718 件の全数について、停止したものが 143 件（8.3\%）、
継続中が 1,575 件である。停止したものの継続期間は中央値 1.62 年、
四分位は 0.64--2.82 年、最大 6.0 年であった。継続中のものは中央値 2.75 年である。
停止したものを事後に集めたのではなく全数から数えているため、
第\ref{sec:survivorship}節の生存条件づけは生じない。

種別ごとの停止率には順序がある。役務 $1/\lambda=11.1$ 年、
高倍率の運用 11.5 年、発行支援 15.3 年、指数 18.1 年、橋渡し 20.1 年、
保険 22.3 年、貸付 22.6 年、利回り 22.7 年である。
\textbf{他の手続きへの依存が薄い種別ほど長い}という向きは読めるが、
予測$\Phi$1 が棄却された以上、この順序を依存の数で説明することはしない。

\paragraph{拘束された状態 $B$。}
表\ref{tab:protocol-quantities}は $B$ について「対応する量を本稿は持たない」と記していた。
区分「内」では手続き上に置かれた資産として直接観測でき、
1,718 件の日次系列が得られる。予測$\Phi$2 で $M(t)$ の代理としたのはこれである。

\paragraph{料率。}
$B$ が測れると、年間の手数料を $B$ で割った料率が定義できる。
ただし\textbf{分母は族によって異なる}。族3（資産の時間貸し）では $B$ が正しい分母だが、
族6（媒介）では取引額が分母であり $B$ は在庫にすぎない。
混ぜると比が発散する。

族3系に限り、$B$ が 100 万ドル以上のもの 144 件で測った料率は次のとおりである。

\begin{center}
\begin{tabularx}{\textwidth}{@{}Xrrr@{}}
\toprule
種別 & $n$ & 中央値 & 四分位 \\
\midrule
裁定 & 7 & 4.97\% & 2.00--9.31\% \\
実物資産 & 20 & 3.85\% & 2.35--5.84\% \\
担保付債務 & 15 & 3.53\% & 0.99--11.86\% \\
流動化した預託 & 19 & 3.12\% & 2.58--4.21\% \\
貸付 & 30 & 2.97\% & 1.77--4.68\% \\
危険の管理 & 21 & 2.83\% & 2.02--4.47\% \\
利回り & 15 & 2.72\% & 0.76--8.01\% \\
利回りの集約 & 14 & 1.69\% & 0.90--4.19\% \\
\bottomrule
\end{tabularx}
\captionof{table}{族3系の料率（年間手数料 $\div$ 拘束された状態）。}
\label{tab:protocol-fee-ladder}
\end{center}

全体の中央値は 3.07\%、四分位は 1.77--5.84\% である。
第\ref{ch:platformfee}章の料率の梯子と同型の並びだが、
\textbf{比較はできない}。第\ref{ch:platformfee}章の料率は取引額に対する率であり、
ここでの料率は残高に対する率である。次元が違う。

\section{本章の限界}
\label{sec:protocol-limits}

\begin{enumerate}[label=(\arabic*)]
\item \textbf{演繹である。} 結論は前提(1)(2)(3)に依存する。
第\ref{sec:protocol-results}節の検定は演繹の含意を測ったものであり、
前提そのものを検証していない。

\item \textbf{判定の対象を絞る規則を要した。} 表\ref{tab:family-assignment-order}は
$\Phi$ を所与としており、台帳の識別子をそのまま当てることはできない。
第\ref{sec:protocol-results}節では判定順序の前に
「これは $\Phi$ か」を置いた（注\ref{rem:is-it-phi}）。
この一段は第\ref{ch:catalog}章の規則に含まれていない。

\item \textbf{両端しか扱っていない。} 注\ref{rem:oracle}のとおり、
区分「外」において第三者の除去の程度は連続的である。
本章は完全に除去された場合と全く除去されない場合のみを扱い、中間を扱わない。
実在する仕組みの多くは中間にある。

\item \textbf{定義の限界が残る。} 第\ref{sec:identifier}節の $\iota$ の観測不能は、
第\ref{sec:phi}章で $N$ を固定したことに由来する。
本章はこれを記録するにとどめ、定義域の拡張は行わない。
第\ref{sec:token-limit}節で述べた請求権の流通と同じ箇所である。

\item \textbf{$\phi$ を扱っていない。} 本章は $\Phi$ の実現可能性のみを論じ、
第\ref{sec:surplus}章の余剰の三分解を適用していない。
手続きが複製可能であることから $\phi_{\rm prod}$ に圧力がかかるという
第\ref{part:unmanned}章の注と同型の議論が成り立ちうるが、
判定には手数料の実測を要するため本章では扱わない。
\end{enumerate}

%============================================================
